\documentclass[11pt]{article}

% ==== Paquetes ====
\usepackage[utf8]{inputenc}
\usepackage[T1]{fontenc}
\usepackage[spanish]{babel}
\usepackage{amsmath, amssymb, amsfonts}
\usepackage{graphicx}
\usepackage{caption}
\usepackage{subcaption}
\usepackage{float}
\usepackage{geometry}
\usepackage{hyperref}
\usepackage{minted}
\usepackage{mdframed}
\geometry{margin=2.3cm}

\usemintedstyle{friendly}
\setminted{fontsize=\scriptsize, breaklines=true}

% ==== Datos de la tarea ====
\title{IEE3864 – Fundamentos de instrumentación astronómica \\ \large Tarea 4}
\author{Nicolás Seguel}
\date{\today}

\begin{document}
\maketitle

\begin{enumerate}
    \item \textbf{Using a glass of n=1.5 design a lens of focal length f=1000 mm}
    
La distancia focal de una lente delgada se puede calcular usando la fórmula de la lente delgada:

\[\frac{1}{f} = (n - 1) \left( \frac{1}{R_1} - \frac{1}{R_2} \right)\]

Para simplificar, podemos asumir que la lente es simétrica, lo que significa que \(R_1 = -R_2\). Esto nos da:

\[\frac{1}{1000} = (1.5 - 1) \left( \frac{1}{R_1} - \frac{1}{-R_1} \right)\]

\[\frac{1}{1000} = 0.5 \left( \frac{2}{R_1} \right)\]

\[\frac{1}{1000} = \frac{1}{R_1}\]

Por lo tanto, \(R_1 = 1000\) mm y \(R_2 = -1000\) mm. Esto significa que la lente es convergente.

    \item \textbf{The lens of the previous problem is used as objective of a telescope, what is the scale on the focal plane of the system in arcsec/mm ?}
    
La escala en el plano focal de un telescopio se puede calcular usando la fórmula:

\[\text{Escala} = \frac{206265}{f}\]

En este caso, \(f = 1000\) mm, por lo que:

\[\text{Escala} = \frac{206265}{1000} = 206.265 \text{ arcsec/mm}\]

    \item \textbf{Using a glass of n=1.5 design a lens of focal length f=-1000 mm}

Utilizando la misma fórmula de la lente delgada. En este caso, queremos que \(f = -1000\) mm y \(n = 1.5\). Asumiendo nuevamente que la lente es simétrica, tenemos:

\[\frac{1}{-1000} = (1.5 - 1) \left( \frac{1}{R_1} - \frac{1}{-R_1} \right)\]
\[\frac{1}{-1000} = 0.5 \left( \frac{2}{R_1} \right)\]
\[\frac{1}{-1000} = \frac{1}{R_1}\]

Por lo tanto, \(R_1 = -1000\) mm y \(R_2 = 1000\) mm. Esto significa que la lente es divergente.

    \item \textbf{Design a lens such that when so=20 cm, si=5 cm. What is the focal length of such a lens?}

Utilizando la fórmula de la lente delgada:

\[\frac{1}{f} = \frac{1}{s_o} + \frac{1}{s_i}\]

Donde \(s_o\) es la distancia del objeto y \(s_i\) es la distancia de la imagen.

\[\frac{1}{f} = \frac{1}{20} + \frac{1}{5} = \frac{1}{4}\]
    

\end{enumerate}

\end{document}